\section{Related Work}
\label{sec:related}

\paragraph{Indic and Devanagari handwritten character recognition.}
The DHCD dataset~\cite{acharya2015dhcd}, introduced in 2015, established the
standard 46-class benchmark for handwritten Devanagari character recognition
and has driven a sustained sequence of architectural improvements.
MallaNet~\cite{malla2025}, the current state-of-the-art on DHCD, couples
residual branch merging with homogeneous filter capsule (HFC) layers to reach
99.71\,\% at 17.32\,M parameters.
Mishra et al.~\cite{mishra2021} independently report 99.72\,\% with a
comparably large architecture.
Across this line of work, each successive generation of models achieves
accuracy improvements measured in hundredths of a percent, while model size
remains large; the trend of ever-larger networks producing ever-smaller gains
motivates a careful examination of whether reported improvements are
statistically distinguishable from measurement noise.

\paragraph{Knowledge distillation and compact model design.}
Hinton et al.~\cite{hinton2015distill} showed that a student network trained
on soft probability outputs (``dark knowledge'') of a larger teacher can match
or approach the teacher's accuracy at substantially lower capacity.
Residual identity mappings~\cite{he2016identity} and
squeeze-and-excitation channel re-weighting~\cite{hu2018senet} have since
become standard building blocks that enable compact architectures to retain
representational power while reducing parameter counts.
Despite the maturity of these techniques, knowledge distillation has rarely
been evaluated on saturated benchmarks where the performance gap between large
and small models is already within the noise of sampling variation; the
question of whether distillation provides a statistically significant
advantage over direct supervised training in this regime remains open.

\paragraph{Statistical significance and benchmark saturation.}
On test sets where top models err on fewer than 0.3\,\% of examples, small
sample counts make raw accuracy rankings unreliable.
McNemar's test~\cite{mcnemar1947} provides the appropriate paired framework:
it conditions on the subset of examples where two models disagree, making it
sensitive to genuine error-pattern differences while being robust to the high
base-rate correct predictions shared by both classifiers.
Wilson score intervals~\cite{wilson1927} give calibrated single-model
accuracy bounds, and calibration analysis~\cite{guo2017calibration}
connects model confidence to empirical accuracy---a prerequisite for
interpreting soft distillation targets.
Applying this suite of tools systematically to a Devanagari recognition
benchmark reveals a methodological gap in prior work: published rankings
have been based on point estimates alone, without paired testing that would
support causal claims about architectural superiority.
\modelname{} fills this gap by reporting exact McNemar $p$-values and Wilson
confidence intervals across every model comparison, demonstrating that the
benchmark is saturated and that statistical parity, rather than raw accuracy
rank, is the appropriate evaluation criterion at this operating point.
